\documentclass[10pt,a4paper]{article}
\usepackage{amsmath,amssymb,amsthm,graphicx}
\title{Low-product-depth iterated matrix multiplication circuit lower bounds over arbitrary fields}
\author{Michael A. Forbes}
\date{doi:10.4230/LIPIcs.CCC.2024.31; CC BY 4.0}
\begin{document}
\maketitle
\noindent\textit{Allowed standard prerequisites (editorial).}\begin{itemize}
\item LST large-characteristic set-multilinearization. Explicitly quoted author Theorem8 citing [9,15]; full statement retained.
\item LST set-multilinear relative-rank bounds with improved BDS parameters. Explicitly quoted author Theorem9 citing [3,15]; complete quoted statement retained.
\end{itemize}\medskip
\section*{Author circuit and polynomial conventions}
Algebraic complexity asks for the best methods to compute multivariate polynomials $f(x_1,\ldots,x_n)$ from the primitive operations of addition and multiplication, using constants from the underlying field $\mathbb F$ for free. These computations are organized as an algebraic circuit, a directed acyclic graph of these primitive operations from the input variables $x_1,\ldots,x_n$. The complexity of this object is primarily governed by its size, which is the number of operations in the circuit. Another complexity measure is the product depth, which is the maximum number of product operations on any input to output path.
The iterated matrix multiplication polynomial $\operatorname{IMM}_{n,d}=(X_1\cdots X_d)_{1,1}$ where $X_i$ is an $n\times n$ symbolic matrix. $F_k$ is the $k$-th Fibonacci number (so $F_0=0$, $F_1=F_2=1$, $F_3=2$, etc.).
\paragraph{Definition 6.} Let $\mathbb F$ a field. Let $\bar x$ be a set of variables with a partition $\bar x=\bar y_1\sqcup\cdots\sqcup\bar y_{d_y}\sqcup\bar z_1\sqcup\cdots\sqcup\bar z_{d_z}$, where $d=d_y+d_z$. Let $Y$ denote the set of all set-multilinear monomials with respect to the partition $\bar y=\bar y_1\sqcup\cdots\sqcup\bar y_{d_y}$, and let $Z$ denote the set of all set-multilinear monomials with respect to the partition $\bar z=\bar z_1\sqcup\cdots\sqcup\bar z_{d_z}$.
Given a polynomial $f(\bar x)$ which is set-multilinear with respect to the above partition, define the coefficient matrix $C_f\in\mathbb F^{Y\times Z}$ by
\[(C_f)_{\bar y^{\bar b},\bar z^{\bar c}}=\operatorname{Coeff}_{\bar y^{\bar b}\bar z^{\bar c}}(f),\]
where $\bar y^{\bar b},\bar z^{\bar c}$ are set-multilinear monomials, and $\operatorname{Coeff}$ takes the coefficient of $\bar y^{\bar b}\bar z^{\bar c}$ in $f$.
The relative rank of $f$ is then defined as
\[\operatorname{relrank}_{\mathbb F}(f):=\frac{\operatorname{rank}_{\mathbb F}(C_f)}{\sqrt{|Y|\cdot|Z|}}.\]
where $\operatorname{rank}_{\mathbb F}(C_f)$ is the matrix rank of $C_f$ over the field $\mathbb F$.
The following set-multilinearization result is a slight extension of what is proven in LST, and follows from their methods as noted by [9].
\paragraph{Definition 7.} Let the variables $\bar x$ be partitioned into $\bar x=\bar x_1\sqcup\cdots\sqcup\bar x_d$. A monomial is set-multilinear (with respect to the partition $\bar x=\bar x_1\sqcup\cdots\sqcup\bar x_d$) if it can be written has $\prod_{i=1}^d(\bar x_i)_{j_i}$ for some $j_1,\ldots,j_d$.
Define the set-multilinear projection (with respect to the partition $\bar x=\bar x_1\sqcup\cdots\sqcup\bar x_d$) to be the linear map $\pi_{\rm sm}:\mathbb F[\bar x]\to\mathbb F[\bar x]$ which is identity on set-multilinear monomials, and zero on all other monomials.
\paragraph{Theorem 8 ([9,15]).} Let $\mathbb F$ be a field of characteristic $\operatorname{char}(\mathbb F)>d$ (or zero), where $d$ is a parameter. Let $f$ be a product-depth $\Delta$ circuit of size $s$. Let the variables $\bar x$ be partitioned into $\bar x=\bar x_1\sqcup\cdots\sqcup\bar x_d$. Then the set-multilinear projection $\pi_{\rm sm}(f)$ has a $(d^d\cdot s)^{O(1)}$-size set-multilinear circuit of product-depth $2\Delta$.
We quote here the summary of the LST lower bound results that we need, incorporating the improved parameters from Bhargav, Dutta, and Saxena [3].
\paragraph{Theorem 9 ([3,15]).} Let $\mathbb F$ a field. Let $\bar x$ be a set of $n$ variables, and $d\le o(\log n)$ a parameter. Then there exists a partition $\bar x=\bar y_1\sqcup\cdots\sqcup\bar y_{d_y}\sqcup\bar z_1\sqcup\cdots\sqcup\bar z_{d_z}$, only depending on $n$ and $d$, where $d=d_y+d_z$, such that any $f(\bar x)$ computed by a set-multilinear circuit of size $s$ and product-depth $\Delta$ has relative rank bounded by
\[\operatorname{relrank}(f)\le s\cdot n^{-\Theta(d^{\frac1{F_{\Delta+2}-1}}/\Delta)}.\]
Further, there exists a set-multilinear polynomial $P$ with $\{0,1\}$-coefficients such that
\[\operatorname{relrank}(P)\ge\frac1{n^{\Theta(1)}},\]
and $P$ can be computed via evaluating the iterated matrix multiplication polynomial $\operatorname{IMM}_{n,d}$ to carefully chosen linear forms.

We now give our characteristic-free version of the above.
\paragraph{Theorem 10.} Let $\mathbb F$ be any field, and $d=o(\log n)$. Then $\operatorname{IMM}_{n,d}$ requires
\[n^{\Theta(d^{\frac1{F_{2\Delta+2}-1}}/\Delta)}\]
size algebraic circuits of product depth $\Delta$.

\section*{Author transfer lemma and its complete proof}
The following is our key lemma that transfers algebraic statements between different fields, in particular showing that an integer matrix having low rank in characteristic zero implies it has low-rank over every field. For our actual needs, we will need to consider matrices with entries that are polynomials with integer coefficients.
\paragraph{Lemma 5.} Let $M\in\mathbb Z[\bar x]^{m\times n}$ be a matrix with integer polynomial entries. Let $\mathbb F$ be any field, and interpret $M\in\mathbb F[\bar x]^{m\times n}$ via the unique non-trivial ring homomorphism $\varphi:\mathbb Z\to\mathbb F$. Then for any $\bar\alpha\in\mathbb F^{\bar x}$,
\[\operatorname{rank}_{\mathbb F}M(\bar\alpha)\le\operatorname{rank}_{\mathbb F(\bar x)}M(\bar x)\le\operatorname{rank}_{\mathbb Q(\bar x)}M(\bar x).\]
\begin{proof}
1) $\operatorname{rank}_{\mathbb F(\bar x)}M(\bar x)\le\operatorname{rank}_{\mathbb Q(\bar x)}M(\bar x)$: In both cases the matrix $M$ is the same, we simply change the field of interpretation. Recall that a matrix is rank $\le s$ iff all $(s+1)\times(s+1)$ submatrices have zero determinant. Let $r=\operatorname{rank}_{\mathbb Q(\bar x)}M(\bar x)$ so that all $(r+1)\times(r+1)$ submatrices of $M[\bar x]$ have zero determinant in $\mathbb Q(\bar x)$. As $M$ has polynomial entries, and the determinant is a polynomial, it follows that under the homomorphism $\varphi$ that these submatrices still have zero determinant. As such, all $(s+1)\times(s+1)$ submatrices of $M$ have zero determinant when viewed as a matrix in $\mathbb F[\bar x]^{m\times n}$, so that $\operatorname{rank}_{\mathbb F(\bar x)}M(\bar x)\le s$.
2) $\operatorname{rank}_{\mathbb F}M(\bar\alpha)\le\operatorname{rank}_{\mathbb F(\bar x)}M(\bar x)$: Let $t=\operatorname{rank}_{\mathbb F(\bar x)}M(\bar x)$, so that all $(t+1)\times(t+1)$ submatrices of $M$ have determinant zero in $\mathbb F[\bar x]$. Define the ring homomorphism $\psi:\mathbb F[\bar x]\to\mathbb F$ by $\bar x\to\bar\alpha$. As above, it then follows that all $(t+1)\times(t+1)$ submatrices of $M$ have zero determinant under this homomorphism. But as $\psi(M(\bar x))=M(\bar\alpha)$, it then follows that all such submatrices of $M(\bar\alpha)$ have zero determinant in $\mathbb F$, so $\operatorname{rank}_{\mathbb F}M(\bar\alpha)\le t$.
\end{proof}
\begin{proof}[Proof of Theorem 10]
Let $\bar x$ be a set of $n$ variables. From Theorem 9, there exists a partition $\bar x=\bar y_1\sqcup\cdots\sqcup\bar y_{d_y}\sqcup\bar z_1\sqcup\cdots\sqcup\bar z_{d_z}$, only depending on $n$ and $d$ (and not on $\mathbb F$), where $d=d_y+d_z$, such that any $f(\bar x)$ computed by a set-multilinear circuit of size $s$ and product-depth $\Delta$ has relative rank bounded by
\[\operatorname{relrank}(f)\le s\cdot n^{-\Theta(d^{\frac1{F_{\Delta+2}-1}}/\Delta)}.\]
Further, there exists a set-multilinear polynomial $P$ with $\{0,1\}$-coefficients such that
\[\operatorname{relrank}(P)\ge\frac1{n^{\Theta(1)}},\]
and $P$ can be computed via evaluating the iterated matrix multiplication polynomial $\operatorname{IMM}_{n,d}$ to carefully chosen linear forms. In particular, any algebraic circuit lower bound for $P$ extends, up to $\operatorname{poly}(n)$ factors, to $\operatorname{IMM}_{n,d}$, so it suffices to prove the lower bound for $P$.
Suppose $P$ (interpreted as a polynomial in $\mathbb F[\bar x]$) is computed by an algebraic circuit $\Phi$ over $\mathbb F$ of size $s$ and product-depth $\Delta$. Create a new algebraic circuit $\Psi$ by replacing each field constant used in $\Phi$ with a distinct variable, so $\Psi$ is size-$s$ product-depth $\Delta$ algebraic circuit over the original variables $\bar x$ along with new variables $\bar w$. Denote $f(\bar x,\bar w)$ to be the polynomial computed by $\Psi$. As $\Psi$ is defined free from field constants, $f$ can be viewed as an integer polynomial $f\in\mathbb Z[\bar x,\bar w]$. Finally, we can relate $P$ and $f$ by undoing the above, so that there are values $\bar\gamma$ from $\mathbb F$ such that $P(\bar x)=f(\bar x,\bar\gamma)$.
Note that $f(\bar x,\bar w)$ may not be set-multilinear, or even of particularly low-degree. However, it does follow that $\pi_{\rm sm}(P(\bar x))=\pi_{\rm sm}(f(\bar x,\bar\gamma))$ as the equality $P(\bar x)=f(\bar x,\bar\gamma)$ is coefficient-wise, so applying $\pi_{\rm sm}$ to each side of equation either keeps the coefficient of $P$ and $f$ the same (if the monomial is set-multilinear) or makes both coefficients zero (if the monomial is not set-multilinear). As $P$ is set-multilinear, we have $P=\pi_{\rm sm}(P)$, so hence $P=\pi_{\rm sm}(f(\bar x,\bar\gamma))$.
Now view $\Psi$ as a circuit with constants over the field $\mathbb Q(\bar w)$, so that $f\in\mathbb Q(\bar w)[\bar x]$. It follows from Theorem 8 that $\pi_{\rm sm}(f)$ has a $(d^d\cdot s)^{O(1)}$-size set-multilinear circuit of product-depth $2\Delta$, and as such,
\[\operatorname{relrank}_{\mathbb Q(\bar w)}(\pi_{\rm sm}(f(\bar x,\bar w)))\le(d^d\cdot s)^{O(1)}\cdot n^{-\Theta(d^{\frac1{F_{2\Delta+2}-1}}/\Delta)}.\]
Using that relative rank is just the (scaled) rank of a matrix, we can invoke Lemma 5, to see that
\[\operatorname{relrank}_{\mathbb F}(\pi_{\rm sm}(f(\bar x,\bar\gamma)))\le\operatorname{relrank}_{\mathbb Q(\bar w)}(\pi_{\rm sm}(f(\bar x,\bar w))).\]
As $P=\pi_{\rm sm}(f(\bar x,\bar\gamma))$, we thus obtain that
\[\frac1{n^{\Theta(1)}}\le\operatorname{relrank}_{\mathbb F}(P)\le(d^d\cdot s)^{O(1)}\cdot n^{-\Theta(d^{\frac1{F_{2\Delta+2}-1}}/\Delta)}\]
which yields the desired lower bound for $s$ (using that $d=o(\log n)$).
\end{proof}

\begin{thebibliography}{99}
\bibitem{source1} Prashanth Amireddy, Ankit Garg, Neeraj Kayal, Chandan Saha, and Bhargav Thankey. Low- depth arithmetic circuit lower bounds: Bypassing set-multilinearization. In Proceedings of the 50th International Colloquium on Automata, Languages and Programming (ICALP 2023), volume 261 of Leibniz International Proceedings in Informatics (LIPIcs), pages 12:1–12:20, 2023. Full version in the Electronic Colloquium on Computational Complexity (ECCC), Technical Report TR22-151. doi:10.4230/LIPICS.ICALP.2023.12.
\bibitem{source2} Robert Andrews. Algebraic hardness versus randomness in low characteristic. In Proceedings of the 35th Annual Computational Complexity Conference (CCC 2020), volume 169 of Leibniz International Proceedings in Informatics (LIPIcs), pages 37:1–37:32, 2020. Full version at arXiv:2005.10885. doi:10.4230/LIPICS.CCC.2020.37.
\bibitem{source3} C. S. Bhargav, Sagnik Dutta, and Nitin Saxena. Improved lower bound, and proof barrier, for constant depth algebraic circuits. In Proceedings of the 47th Internationl Symposium on the Mathematical Foundations of Computer Science (MFCS 2022), volume 241 of Leibniz International Proceedings in Informatics (LIPIcs), pages 18:1–18:16, 2022. doi:10.4230/ LIPICS.MFCS.2022.18.
\bibitem{source4} Radu Curticapean, Nutan Limaye, and Srikanth Srinivasan. On the VNP-hardness of some monomial symmetric polynomials. In 42nd International Conference on Foundations of Software Technology and Theoretical Computer Science (FSTTCS 2022), volume 250 of Leibniz International Proceedings in Informatics (LIPIcs), pages 16:1–16:14, 2022. Full version in the Electronic Colloquium on Computational Complexity (ECCC), Technical Report TR22-139. doi:10.4230/LIPICS.FSTTCS.2022.16.
\bibitem{source5} Ismor Fischer. Sums of like powers of multivariate linear forms. Mathematics Magazine, 67(1):59–61, 1994. URL: http://www.jstor.org/stable/2690560.
\bibitem{source6} Philippe Flajolet and Robert Sedgewick. Analytic combinatorics. Cambridge University Press, 2009. doi:10.1017/CBO9780511801655.
\bibitem{source7} Michael A. Forbes, Amir Shpilka, and Ben Lee Volk. Succinct hitting sets and barriers to proving algebraic circuits lower bounds. In Proceedings of the 49th Annual ACM Symposium on Theory of Computing (STOC 2017), pages 653–664, 2017. Full version at arXiv:1701.05328. doi:10.1145/3055399.3055496.
\bibitem{source8} Hervé Fournier, Nutan Limaye, Srikanth Srinivasan, and Sébastien Tavenas. On the power of homogeneous algebraic formulas. Electronic Colloquium on Computational Complexity (ECCC), TR23-191, 2023. URL: https://eccc.weizmann.ac.il/report/2023/191.
\bibitem{source9} Nashlen Govindasamy, Tuomas Hakoniemi, and Iddo Tzameret. Simple hard instances for low-depth algebraic proofs. In Preliminary version in the 63rd Annual IEEE Symposium on Foundations of Computer Science (FOCS 2022), pages 188–199, 2022. Full version at arXiv:2205.07175. doi:10.1109/FOCS54457.2022.00025.
\bibitem{source10} Joshua A. Grochow, Mrinal Kumar, Michael E. Saks, and Shubhangi Saraf. Towards an algebraic natural proofs barrier via polynomial identity testing. arXiv, 1701.01717, 2017. URL: http://arxiv.org/abs/1701.01717.
\bibitem{source11} Joshua A. Grochow and Toniann Pitassi. Circuit complexity, proof complexity, and polynomial identity testing. In Proceedings of the 55th Annual IEEE Symposium on Foundations of Computer Science (FOCS 2014), pages 110–119, 2014. Full version at arXiv:1404.3820. doi:10.1109/FOCS.2014.20.
\bibitem{source12} Alan Guo, Swastik Kopparty, and Madhu Sudan. New affine-invariant codes from lifting. In Proceedings of Innovations in Theoretical Computer Science (ITCS 2013), pages 529–540, 2013. Full version at arXiv:1208.5413. doi:10.1145/2422436.2422494.
\bibitem{source13} Ankit Gupta, Pritish Kamath, Neeraj Kayal, and Ramprasad Saptharishi. Arithmetic circuits: A chasm at depth three. In Proceedings of the 54th Annual IEEE Symposium on Foundations of Computer Science (FOCS 2013), pages 578–587, 2013. Full version in the Electronic Colloquium on Computational Complexity (ECCC), Technical Report TR13-026. doi:10. 1109/FOCS.2013.68.
\bibitem{source14} Erich L. Kaltofen. Factorization of polynomials given by straight-line programs. In Silvio Micali, editor, Randomness and Computation, volume 5 of Advances in Computing Research, pages 375–412. JAI Press, Inc., Greenwich, CT, USA, 1989. URL: http://www.math.ncsu. edu/~kaltofen/bibliography/89/Ka89\_slpfac.pdf.
\bibitem{source15} Nutan Limaye, Srikanth Srinivasan, and Sébastien Tavenas. Superpolynomial lower bounds against low-depth algebraic circuits. In Preliminary version in the 62nd Annual IEEE Sym- posium on Foundations of Computer Science (FOCS 2021), pages 804–814, 2022. Full version in the Electronic Colloquium on Computational Complexity (ECCC), Technical Report TR21-081. doi:10.1109/FOCS52979.2021.00083.
\bibitem{source16} Nutan Limaye, Srikanth Srinivasan, and Sébastien Tavenas. Guest column: Lower bounds against constant-depth algebraic circuits. SIGACT News, 53(2):40–62, 2022. doi:10.1145/ 3544979.3544989.
\bibitem{source17} Henryk Minc. Evaluation of permanents. Proc. Edinburgh Math. Soc. (2), 22(1):27–32, 1979. doi:10.1017/S0013091500027760.
\bibitem{source18} Ran Raz. Elusive functions and lower bounds for arithmetic circuits. Theory of Computing, 6(1):135–177, 2010. 40th Annual ACM Symposium on Theory of Computing (STOC 2008). doi:10.4086/TOC.2010.V006A007.
\bibitem{source19} Alexander A. Razborov. Lower bounds on the size of bounded depth circuits over a complete basis with logical addition. Matematicheskie Zametki, 41(4):598–607, April 1987.
\bibitem{source20} Amir Shpilka and Avi Wigderson. Depth-3 arithmetic circuits over fields of characteristic zero. Computational Complexity, 10(1):1–27, 2001. Preliminary version in the 14th Annual IEEE Conference on Computational Complexity (CCC 1999). doi:10.1007/PL00001609.
\bibitem{source21} Roman Smolensky. Algebraic methods in the theory of lower bounds for boolean circuit complexity. In Proceedings of the 19th Annual ACM Symposium on Theory of Computing (STOC 1987), pages 77–82, 1987. doi:10.1145/28395.28404.
\bibitem{source22} Ryan Williams. Finding paths of length k in O∗ (2k ) time. Inf. Process. Lett., 109(6):315–318, 2009. doi:10.1016/J.IPL.2008.11.004.
\end{thebibliography}
\end{document}
